\section{Smoothing and reverse comparison}\label{sec:reverse}

Proposition~\ref{prop:proposal} controls ${\rm KL}(Q_l\Vert P)$.
For coverage we need the opposite direction: a set of appreciable
Gibbs mass must receive a quantitative amount of proposal mass.
Conditioning on the spin makes the Gibbs disorder law strongly
log-concave. A small Gaussian smoothing then turns the forward
comparison into the required reverse bound.

For the remainder of the proof fix
\[
 f=1/100,\qquad d=f/100,\qquad b=d/100,
 \qquad \delta=n^{-f/4}.
\]
The longer proposal horizon $l=f$ will provide coverage of sets with
Gibbs mass at least $n^{-d}$. We reserve the shorter horizon $l=b$
for Section~\ref{sec:connections}, where its stronger exceptional-event
estimate will control a sequence of proposals. For any joint law of
disorder and spin, define $\mathcal S$ by retaining the spin
and replacing its disorder by
\[
 G^\circ=\sqrt{1-\delta}\,G+\sqrt\delta\,\bar G,
\]
where $\bar G$ is an independent standard Gaussian vector.
Both $\mathcal SP$ and $\mathcal SQ_f$ have standard Gaussian
disorder marginal. Denote the conditional spin kernel of
$\mathcal SQ_f$ by $\nu(\cdot\mid G^\circ)$. Gaussian regression
gives a useful way to generate it: given $G^\circ$, run the
$l=f$ proposal on
$\sqrt{1-\delta}\,G^\circ+\sqrt\delta\,G'$, using fresh independent
standard Gaussian $G'$ and independent auxiliary inputs.

\subsection{Why smoothing reverses the entropy comparison}

We prove
\[
 {\rm KL}(\mathcal SP\Vert\mathcal SQ_f)
 \le\delta^{-1}{\rm KL}(Q_f\Vert P).
 \tag{8}\label{eq:8}
\]
The two unsmoothed laws have the same uniform spin marginal. Given
$x$, the conditional disorder density under $P$ has negative logarithm,
up to a normalizing constant,
\[
 U_x(G)=\tfrac12\|G\|^2-\beta\widehat H_G(x)+F(G).
\]
Its Hessian is at least the identity: the Hamiltonian is linear in
$G$, while $F$ is a convex log-sum-exp of linear functions.

We need the following consequence of this convexity. A law $\mu$
with density proportional to $e^{-U}$, where $\nabla^2U\ge I$,
and a law $\lambda\ll\mu$ admit couplings whose mean squared
displacement is at most $2{\rm KL}(\lambda\Vert\mu)$, or tends
to that bound. This is the quadratic transport inequality; its
triangular-transport form appears in
\cite[Corollary~3.10]{BKM2005}. We recall the proof for the
finite-dimensional densities used here.

First restrict to a box and suppose both densities are smooth and
strictly positive up to its boundary. Map $\mu$ to $\lambda$ by
successive increasing conditional quantile maps. The resulting
triangular map $T$ has positive diagonal derivatives, whose product
is its Jacobian determinant. The entropy formula is
\[
 {\rm KL}(\lambda\Vert\mu)
 =\int\bigl(U(T(z))-U(z)-\log\det\nabla T(z)\bigr)\,d\mu(z).
\]
Strong convexity bounds the first difference below by
$\nabla U(z)\cdot(T(z)-z)+\tfrac12\|T(z)-z\|^2$.
Integration by parts replaces the linear term by
$\operatorname{div}(T-\mathrm{id})$. There is no boundary term:
on a face where the $i$th coordinate is an endpoint of the box,
its conditional quantile map fixes that endpoint, so the normal
component $T_i-z_i$ vanishes. Tangential motion does not matter.
The remaining Jacobian terms are nonnegative because
$z-1-\log z\ge0$ for every positive diagonal derivative $z$.
This proves the coupling bound on the box.

The densities needed here satisfy the approximation requirements.
For fixed $n$, the conditional proposal density is bounded by a
constant times Gaussian density, since its spin probability is at
most one. Truncate both conditional laws to increasing boxes and
normalize them. On each box approximate $\lambda$ by smoothing and
mixing with a small positive constant density. Its density is bounded,
and the smooth reference density is bounded away from zero on the
box, so their relative entropies with respect to the reference converge.
The original laws have Gaussian tails up to constants for fixed $n$,
and finite relative entropy.
Thus the truncated entropies converge; tightness and lower
semicontinuity pass the coupling bound to the original laws.

Apply this result with $\mu=P(\cdot\mid x)$ and
$\lambda=Q_f(\cdot\mid x)$. For any coupling of two disorder
values $g,\widetilde g$, the Gaussian smoothing kernels satisfy
\[
 {\rm KL}\bigl(N(\sqrt{1-\delta}\,g,\delta I)
       \Vert N(\sqrt{1-\delta}\,\widetilde g,\delta I)\bigr)
 =\frac{1-\delta}{2\delta}\|g-\widetilde g\|^2.
\]
The convolution comparison is also recorded in
\cite[Lemma~11]{LugosiNeu2022}. Joint convexity of entropy bounds
the divergence between the mixtures by this averaged kernel
quantity. Average over the common spin marginal and use the entropy
chain rule. The transport bound then proves \eqref{eq:8}.

\subsection{Coverage conditional on a target and its base}

The smoothed entropy comparison is still a joint-law statement. We
now convert it to a statement for the original Gibbs law at a
realized target disorder. The following proposition records both the
entropy estimate and the coverage consequence used in the next section.

\begin{proposition}[Coverage at a target disorder]\label{prop:coverage}
Let $G,\bar G$ be independent standard Gaussian disorder vectors, set
$G^\circ=\sqrt{1-\delta}\,G+\sqrt\delta\,\bar G$, and let
$\nu(\cdot\mid G^\circ)$ be the conditional spin kernel of
$\mathcal SQ_f$ defined above. Then
\[
 \mathbb E\,{\rm KL}\bigl(\pi_G\Vert\nu(\cdot\mid G^\circ)\bigr)
 \le O_\beta(n\sqrt\delta)+2n^{1-f/2}/\delta.
 \tag{9}\label{eq:9}
\]
With probability tending to one over this target/base pair,
\[
 {\rm KL}\bigl(\pi_G\Vert\nu(\cdot\mid G^\circ)\bigr)
 \le n^{1-4d}.
 \tag{10}\label{eq:10}
\]
On the same event, simultaneously for every set
$A\subseteq\{-1,1\}^n$ with $\pi_G(A)\ge n^{-d}$,
\[
 \nu(A\mid G^\circ)\ge\exp(-n^{1-2d}).
\]
The set $A$ may be chosen after the target and base are realized.
\end{proposition}

\begin{proof}
Let $\overline\pi(\cdot\mid G^\circ)$ be the conditional spin
law of $\mathcal SP$, namely the conditional mixture of $\pi_G$.
Inserting this law into the logarithmic ratio gives the exact
decomposition
\[
 \begin{aligned}
 \mathbb E\,{\rm KL}\bigl(\pi_G\Vert\nu(\cdot\mid G^\circ)\bigr)
 &=\mathbb E\,{\rm KL}\bigl(\pi_G\Vert
                 \overline\pi(\cdot\mid G^\circ)\bigr)\\
 &\quad+{\rm KL}(\mathcal SP\Vert\mathcal SQ_f).
 \end{aligned}
\]
Both expectations are over the target/base pair in the statement.
The conditional mixture minimizes averaged divergence among spin laws
depending only on the base. The first term on the right is therefore
at most $\mathbb E{\rm KL}(\pi_G\Vert\pi_{G^\circ})$.

For any two Hamiltonians the Gibbs divergence is bounded by twice
$\beta$ times the supremum of their absolute energy difference.
Here those differences, for each fixed spin, are centered Gaussians
with variance $O(n\delta)$. A union bound, or exponential moments
over the $2^n$ spins, bounds the expected supremum by
$O(n\sqrt\delta)$. Combining this estimate with \eqref{eq:8}
and \eqref{eq:5} proves \eqref{eq:9}.

Since $4d<f/8$ and $4d<f/4$, Markov's inequality applied to
\eqref{eq:9} gives \eqref{eq:10} with probability tending to one.
On this single event, let $A$ be any set with $\pi_G(A)\ge n^{-d}$.
For $s=\pi_G(A)$ and $t=\nu(A\mid G^\circ)$, entropy
contraction to the two-point partition gives
$s\log(1/t)\le {\rm KL}(\pi_G\Vert\nu)+\log2$.
Hence
\[
 \log(1/t)\le n^d(n^{1-4d}+\log2)\le n^{1-2d}
\]
for sufficiently large $n$. The implication is deterministic once
the target and base are fixed, which proves the simultaneous assertion.
\end{proof}
